\section{Случайни величини. Свойства}

\subsection{Въведение и интуиция}
Често при разглеждане на вероятностни модели, множеството от елементарни събития $\Omega$ може да съдържа обекти с много сложна структура — например хора, молекули, траектории на брауново движение или шесторки числа в тотото. Едно такова множество може да бъде характеризирано чрез всевъзможни характеристики (височина, тегло, доходи и др.).

На практика обаче, за даден експеримент или изследване най-често се интересуваме от конкретна характеристика, която по възможност е числова. Искаме на всяко елементарно събитие $\omega \in \Omega$ да съпоставим реално число. Това класическо изображение наричаме случайна величина. То ни позволява да забравим за сложната структура на елементарните събития и да работим изцяло с техните числови характеристики, което за повечето експерименти е абсолютно достатъчно.

\subsection{Формална дефиниция}
Нека $V = (\Omega, \mathcal{A}, \mathbb{P})$ е вероятностно пространство.
Разглеждаме изображение $X: \Omega \rightarrow \mathbb{R}$.

Ако имаме интервал (или произволно множество) $I \subset \mathbb{R}$, под пълния първообраз на $I$ чрез $X$ ще разбираме множеството от онези елементарни събития, чиято числова характеристика попада в $I$:
\[ X^{-1}(I) = \{ \omega \in \Omega : X(\omega) \in I \} \]
Например, ако $I = (a, b)$, множеството е:
\[ X^{-1}((a, b)) = \{ \omega \in \Omega : a < X(\omega) < b \} \]
което за краткост ще записваме просто като $\{ a < X < b \}$, като винаги подразбираме, че това е подмножество на $\Omega$.

За да можем да оперираме математически строго, ни е необходима специална структура. Трябва да можем да придаваме вероятности на множествата, които ни интересуват.

\begin{keydefn}
Изображението $X: \Omega \rightarrow \mathbb{R}$ се нарича \emph{случайна величина}, ако за всяко реално число $b \in \mathbb{R}$, множеството
\[ X^{-1}((-\infty, b)) = \{ \omega \in \Omega : X(\omega) < b \} \]
принадлежи на $\sigma$-алгебрата $\mathcal{A}$. Кратко записваме $\{ X < b \} \in \mathcal{A}$.
\end{keydefn}

Това условие означава, че можем да придаваме вероятност на такива множества — можем да изчисляваме вероятността $\mathbb{P}(X < b)$. Тоест, можем да отговорим на въпроси като „каква е вероятността доходът на човек да е по-малък от 1000 лева?“ или „каква е вероятността да тежи по-малко от 70 кг?“.

\begin{cor}
Ако $X$ е случайна величина, то за всеки две реални числа $a < b$, множествата $\{ a < X < b \}$, $\{ a \le X \le b \}$, $\{ a \le X < b \}$ и т.н. принадлежат на $\sigma$-алгебрата $\mathcal{A}$. Следователно ние можем да измерваме и да изчисляваме вероятностите случайната величина да попада в произволен интервал, бил той отворен, затворен, полуотворен или безкраен.
\end{cor}

\subsection{Операции със случайни величини}
Нека $V = (\Omega, \mathcal{A}, \mathbb{P})$ е вероятностно пространство и нека $X$ и $Y$ са две случайни величини, дефинирани върху него. Тук е изключително важно двете случайни величини да „вървят“ с едно и също вероятностно пространство, за да можем да дефинираме операциите събиране и умножение върху едни и същи елементарни събития. Тогава са изпълнени следните свойства:
\begin{enumerate}
    \item $c X$ е случайна величина за всяко $c \in \mathbb{R}$.
    \item $X + Y$ и $X - Y$ са случайни величини в $V$.
    \item $X Y$ е случайна величина.
    \item $\frac{X}{Y}$ е случайна величина, при условие че събитието $\{Y = 0\}$ е с вероятност 0. Тоест, изображението не връща стойност 0 (например никой няма 0 кг тегло) и можем да делим на него.
\end{enumerate}

\textbf{Скица на доказателство за $c X$:}
Ще разгледаме случая, когато константата $c < 0$. Нека $Z = c X$. Искаме да проверим, че за всяко реално число $b$, събитието $\{ Z < b \} \in \mathcal{A}$:
\[ \{ c X < b \} = \left\{ X > \frac{b}{c} \right\} \]
Тъй като $c < 0$, знакът на неравенството се обръща. Множеството $\left\{ X > \frac{b}{c} \right\}$ е точно $X^{-1}\left(\left(\frac{b}{c}, \infty\right)\right)$, което въз основа на предходното следствие принадлежи на $\sigma$-алгебрата $\mathcal{A}$. (Ако $c > 0$, знакът се запазва и доказателството следва директно от дефиницията).

\textbf{Идея на доказателството за $X + Y$:}
За да покажем, че сумата $X+Y$ е случайна величина, трябва да проверим дали множеството $\{ X + Y < b \}$ е в $\mathcal{A}$. Не е трудно (въпреки че не е съвсем тривиално) да се покаже, че това събитие може да се представи като изброимо обединение по всички рационални числа $q \in \mathbb{Q}$:
\[ \{ X + Y < b \} = \bigcup_{q \in \mathbb{Q}} ( \{ X < q \} \cap \{ Y < b - q \} ) \]
Ако $\omega$ принадлежи на това множество, то сумата $X(\omega) + Y(\omega) < b$. Виждаме, че $\{ X < q \} \in \mathcal{A}$ (по дефиницията за $X$) и $\{ Y < b - q \} \in \mathcal{A}$ (по дефиницията за $Y$). Тъй като $\mathcal{A}$ е $\sigma$-алгебра, сечението на две нейни множества отново е в $\mathcal{A}$. Тъй като множеството на рационалните числа $\mathbb{Q}$ е изброимо, изброимо обединение на елементи от $\mathcal{A}$ също принадлежи на $\mathcal{A}$. С това сложният въпрос се свежда до работа с отделните случайни величини.

\section{Индикаторни функции и техните свойства}

Преди да преминем към конкретната дефиниция на дискретни случайни величини, е полезно да въведем едно понятие, което не е директно свързано с вероятностите, а по-скоро с множества.

\begin{defn}
Нека $A$ е подмножество на $\Omega$. Индикаторната функция $\ind_A : \Omega \rightarrow \{0, 1\}$ се дефинира по следния начин:
\[
\ind_A(\omega) =
\begin{cases}
1, & \text{ако } \omega \in A \\
0, & \text{ако } \omega \notin A
\end{cases}
\]
\end{defn}

Това е най-простият прототип за бинарен експеримент (успех/неуспех). Цялото сложно пространство $\Omega$ се разцепва на две части: интересуваме се само дали събитието се е случило (стойност 1) или не (стойност 0). Използването на индикаторни функции е изключително удобно, защото ни позволява алгебрично да оперираме със сечения и обединения.

\textbf{Свойства на индикаторните функции:}
За произволни множества $A, B \subset \Omega$:
\begin{enumerate}
    \item \textbf{Сечение:} $\ind_{A \cap B} = \ind_A \cdot \ind_B$. \\
    Това равенство е вярно, защото $\ind_{A \cap B}(\omega) = 1$ тогава и само тогава, когато едновременно $\ind_A(\omega) = 1$ и $\ind_B(\omega) = 1$. Ако дори едното е 0, произведението пропада в 0.
    \item \textbf{Обединение:} $\ind_{A \cup B} = \ind_A + \ind_B - \ind_{A \cap B} = \ind_A + \ind_B - \ind_A \cdot \ind_B$. \\
    Тази формула директно отразява принципа на включване и изключване за две множества.
    \item \textbf{Допълнение:} $\ind_{A^c} = 1 - \ind_A$. \\
    Ако индикаторът за $A$ е 1, то $1 - 1 = 0$ (не е в допълнението), и обратно.
\end{enumerate}

С помощта на индикаторите, известни закони (като законите на Де Морган) се доказват алгебрично, без чертаене на диаграми на Вен:
\begin{align*}
\ind_{(A \cup B)^c} &= 1 - \ind_{A \cup B} = 1 - (\ind_A + \ind_B - \ind_A \ind_B) = 1 - \ind_A - \ind_B + \ind_A \ind_B \\
 &= (1 - \ind_A)(1 - \ind_B) = \ind_{A^c} \ind_{B^c} = \ind_{A^c \cap B^c}
\end{align*}
От равенството на индикаторните функции следва равенството на самите множества: $(A \cup B)^c = A^c \cap B^c$.

\begin{prop}
Ако множеството $A \in \mathcal{A}$, то индикаторната функция $\ind_A$ е случайна величина.
\end{prop}

\begin{proof}
Разглеждаме множеството $\{ \ind_A < b \}$ за произволно реално $b$:
\[
\{ \ind_A < b \} =
\begin{cases}
\emptyset, & \text{ако } b \le 0 \\
A^c, & \text{ако } 0 < b \le 1 \\
\Omega, & \text{ако } b > 1
\end{cases}
\]
Първият случай е така, защото индикаторът е винаги $\ge 0$; третият — защото е винаги $\le 1$; в средния случай единствената възможност е индикаторът да е $0$, т.е. $\omega \in A^c$.

Понеже $A \in \mathcal{A}$, неговото допълнение $A^c \in \mathcal{A}$. Празната съвкупност $\emptyset$ и цялото множество $\Omega$ също принадлежат на $\mathcal{A}$. Следователно и в трите възможни случая множеството попада в $\sigma$-алгебрата. Ето защо $\ind_A$ е случайна величина.
\end{proof}

Това е фундаментът („тухличката“), върху който се изгражда теорията на мярката и интеграла, работейки с линейни комбинации на такива индикатори.

\section{Дискретни случайни величини}\label{sec:discrete-random-variables}

\subsection{Дефиниция и разпределение}
Преминаваме към най-реалистичния клас от случайни величини — дискретните. В крайна сметка, светът и компютрите оперират дискретно и ние не можем да съхраняваме и оперираме директно с неизброима безкрайност.

Нека $V = (\Omega, \mathcal{A}, \mathbb{P})$ е вероятностно пространство. Нека разполагаме с редица от различни реални числа $x_j$ (която може да бъде крайна $x_1, \dots, x_n$ или изброимо безкрайна $x_1, x_2, \dots$) и съответна пълна група от събития $H_j \in \mathcal{A}$. (Пълна група означава, че тези събития са непресичащи се и обединението им изчерпва цялото $\Omega$).

\begin{defn}
Случайната величина $X$, дефинирана чрез
\[ X = \sum_{j} x_j \ind_{H_j} \]
се нарича \emph{дискретна случайна величина}.
\end{defn}

Тъй като събитията $H_j$ образуват пълна група, за всяко $\omega \in \Omega$ само един индикатор от тази сума не е нула. Следователно $X(\omega)$ винаги приема една конкретна стойност $x_j$. Самата величина $X$ е гарантирано случайна величина, тъй като (както видяхме) индикаторите са случайни величини, а тяхна линейна комбинация също е случайна величина.

\begin{defn}
Под \emph{разпределение} на дискретната случайна величина $X$ разбираме таблицата, съдържаща всички възможни стойности на $X$ и техните съответни вероятности $p_j$:
\[
\begin{array}{c|c|c|c|c}
X & x_1 & x_2 & \dots & x_n \\
\hline
\mathbb{P} & p_1 & p_2 & \dots & p_n
\end{array}
\]
Тук $p_j$ означава вероятността $X$ да приеме стойност $x_j$, тоест вероятността да настъпи събитието $H_j$ (записвано още като $\{X = x_j\}$). В сила са условията $p_j > 0$ и $\sum_j p_j = 1$.
\end{defn}

С тази таблица нашият модел напълно се абстрахира от сложното вероятностно пространство $\Omega$. Вече не се интересуваме дали зад модела стоят мъже и жени, температури или тегла. Важното е само какви стойности приема $X$ и с какви вероятности.

\subsection{Равенство по разпределение}
Много важно понятие в вероятностите е равенството по разпределение. Нека $X$ и $Y$ са дискретни случайни величини (които могат да идват от съвсем различни експерименти и вероятностни пространства).

\begin{defn}
Казваме, че $X$ е \emph{равна по разпределение} на $Y$ (записваме $X \stackrel{d}{=} Y$, от distribution), ако техните таблици на разпределение съвпадат. Тоест, за всяка възможна стойност $x_j$ е изпълнено:
\[ \mathbb{P}(X = x_j) = \mathbb{P}(Y = x_j) \]
\end{defn}

\begin{example}[Живот на процесор]\label{ex:04-1}
Нека $X$ измерва живота на даден процесор в дни. $X$ може да приема стойности $0, 1, 2, \dots$ (безкрайна редица) с вероятности $p_0, p_1, p_2, \dots$.
В друг модел, нека фирмата има правило да сменя процесорите най-много на 4000-ния ден. Тогава новата случайна величина $Y$ ще приема стойности от 0 до 4000. Нейните вероятности за $k < 4000$ съвпадат с тези на $X$, но за стойността 4000 вероятността акумулира всички останали: $\mathbb{P}(Y=4000) = \sum_{j=4000}^{\infty} p_j$. Следователно, въпреки че описват сходен процес, $X$ и $Y$ нямат еднакви таблици и не са равни по разпределение.

\end{example}

\begin{example}[Монета и зар]\label{ex:04-2}
Нека $X$ отразява хвърлянето на честна монета ($X=1$ за ези, $X=0$ за тура; и двете с вероятност $1/2$). Нека $Y$ отразява хвърлянето на честен зар, като отчитаме само дали числото е четно ($Y=1$) или нечетно ($Y=0$). Вероятностите за четно и нечетно при зар отново са $1/2$. Следователно таблиците им на разпределение са идентични и $X \stackrel{d}{=} Y$, въпреки че физическите генератори на случайност (монета и зар) са коренно различни.

\end{example}

\subsection{Хистограма}
За нагледно представяне на разпределението често се използва хистограма. Тя се построява като по хоризонталната ос се отбелязват стойностите $x_j$, а по вертикалната се издигат стълбчета с височина $p_j$.

\begin{example}\label{ex:04-3}
Да разгледаме хвърлянето на два честни зара, които се хвърлят независимо. Интересуваме се от сумата им. Възможните стойности са числата от 2 до 12. Вероятностите за всяка сума са съответно:
\[
\begin{array}{c|c|c|c|c|c|c|c|c|c|c|c}
\text{Стойност} & 2 & 3 & 4 & 5 & 6 & 7 & 8 & 9 & 10 & 11 & 12 \\
\hline
\text{Вероятност} & \frac{1}{36} & \frac{2}{36} & \frac{3}{36} & \frac{4}{36} & \frac{5}{36} & \frac{6}{36} & \frac{5}{36} & \frac{4}{36} & \frac{3}{36} & \frac{2}{36} & \frac{1}{36}
\end{array}
\]
Хистограмата за този експеримент ще изглежда като симетрична триъгълна пирамида, с най-висок връх (вероятност $6/36$) при стойността 7.

\end{example}

\section{Смяна на променливите на дискретни случайни величини}

Понякога се налага първоначалната случайна величина да бъде трансформирана в друга, по-удобна за нас. Например, ако разполагаме с резултат от зарче (1 до 6), но се интересуваме само дали е четен или не.

\subsection{Функция на една дискретна случайна величина}
Нека $X$ е дискретна случайна величина, а $g: \mathbb{R} \rightarrow \mathbb{R}$ е функция, която е добре дефинирана поне върху възможните стойности на $X$. Тогава $Y = g(X)$ също е дискретна случайна величина:
\[ Y = g(X) = \sum_{j} g(x_j) \ind_{H_j} \]

Ако функцията $g$ изобразява няколко различни стойности $x_i, x_j$ в едно и също число (както при четните числа от зара, които отиват в 1), ние можем да преномерираме и да обединим съответните колонки в таблицата на разпределението, като сумираме техните вероятности:
\[ \mathbb{P}(Y = y) = \sum_{j : g(x_j) = y} \mathbb{P}(X = x_j) \]

\begin{example}\label{ex:04-4}
Да се върнем към модела с процесора, чийто живот е $X \in \{0, 1, 2, \dots\}$. Искаме да знаем само дали процесорът е бил фабрично дефектен в нулевия ден. Дефинираме функцията $g(0) = 0$ (дефектен) и $g(x) = 1$ за $x \ge 1$ (работещ поне 1 ден). Тогава $Y = g(X)$ ще приема само две стойности:
\[ Y = 0 \cdot \ind_{\{X=0\}} + 1 \cdot \ind_{\{X \ge 1\}} \]
Табличката на разпределение на $Y$ става изключително проста: стойност $0$ с вероятност $p_0$ и стойност $1$ с вероятност $1 - p_0$. Така си спестяваме дългото и сложно разпределение на $X$, фокусирайки се само върху дефектите.

\end{example}

\subsection{Функция на две дискретни случайни величини}
Когато сменяме променливите на две случайни величини едновременно, е ключово те да са дефинирани върху \emph{едно и също} вероятностно пространство $V$.

Нека $X$ и $Y$ са две дискретни случайни величини във $V$, със съответните си пълни групи събития $H_j$ и $H_k^*$. Нека $g: \mathbb{R} \times \mathbb{R} \rightarrow \mathbb{R}$ е функция на две променливи (например $g(x,y) = x+y$). Тогава случайната величина $Z = g(X, Y)$ е:
\[ Z = g(X, Y) = \sum_{j, k} g(x_j, y_k) \ind_{H_j} \ind_{H_k^*} = \sum_{j, k} g(x_j, y_k) \ind_{H_j \cap H_k^*} \]
Тук използваме свойството на индикаторните функции, че произведението им е равно на индикатора на тяхното сечение. Това сечение $H_j \cap H_k^*$ отговаря на едновременното настъпване на събитието $X$ да приеме стойност $x_j$ и $Y$ да приеме стойност $y_k$.

\section{Независимост на дискретни случайни величини}

Концепцията за независимост е ключова. Ако имаме две дискретни случайни величини $X$ и $Y$ във $V$, интуитивно те са независими, ако стойността, която приема едната, не носи никаква информация за стойността на другата.

\begin{defn}[независимост на две случайни величини]
$X$ и $Y$ са \emph{независими} ($X \perp Y$), тогава и само тогава, когато за всеки две възможни техни стойности $x_i$ и $y_j$ е изпълнено:
\[ \mathbb{P}(X = x_i, Y = y_j) = \mathbb{P}(X = x_i) \mathbb{P}(Y = y_j) \]
Тоест, вероятността двете случайни величини едновременно да приемат конкретни стойности (настъпването на сечението на събитията) се разпада като произведение от техните индивидуални вероятности. Това е директен аналог на независимостта на събития: $\mathbb{P}(A_i \cap B_j) = \mathbb{P}(A_i)\mathbb{P}(B_j)$. За да направим това сечение, отново ни трябва $X$ и $Y$ да са в едно и също вероятностно пространство.
\end{defn}

\begin{defn}[независимост в съвкупност]
Нека имаме $n$ дискретни случайни величини $X_1, X_2, \dots, X_n$. Те се наричат независими в съвкупност, тогава и само тогава, когато за всяко подмножество от техни индекси $\{j_1, j_2, \dots, j_m\} \subset \{1, \dots, n\}$ (с размер $m \ge 2$) и за всички техни възможни стойности $x_{j_1}, \dots, x_{j_m}$ е вярно:
\[ \mathbb{P}(X_{j_1} = x_{j_1}, X_{j_2} = x_{j_2}, \dots, X_{j_m} = x_{j_m}) = \mathbb{P}(X_{j_1} = x_{j_1}) \mathbb{P}(X_{j_2} = x_{j_2}) \dots \mathbb{P}(X_{j_m} = x_{j_m}) \]
Записът може да изглежда сложен, но идеята му е проста: каквото и подмножество от случайните величини да вземете, съвместната вероятност за всяка тяхна конфигурация се представя като произведение от маргиналните (индивидуалните) вероятности на всяко едно от събитията.
\end{defn}

\section{Задачи}

\textbf{Домашна работа.} Нека $X$ и $Y$ са случайни величини върху едно и също вероятностно пространство. Довършете доказателството на равенството
\[
\{X+Y<b\}
=\bigcup_{q\in\mathbb{Q}}\bigl(\{X<q\}\cap\{Y<b-q\}\bigr),
\]
като докажете включването отляво надясно.
